\begin{table}[H]
\floatleft
\caption{Issue-domain Susceptibility Summary}
\label{tab:domains}
\vspace{2pt}
\begingroup\footnotesize
\begin{tabularx}{\linewidth}{@{}Y*{7}{c}@{}}
\toprule
Issue domain & Leaves & $n$ & Mdn & Mean & IQR (Q1, Q3) & \% toward goal & $|r_{\mathrm{rb}}|$ vs ref \\
\midrule
Critical Infrastructure and Energy Sovereignty & 12 & 1{,}429 & 19.0 & 22.5 & (14.3, 26.2) & 94.4 & 0.46 \\
Supranational and Regional Integration & 7 & 1{,}428 & 16.8 & 20.6 & (7.3, 27.4) & 81.2 & 0.22 \\
Democratic Resilience and Institutions & 25 & 1{,}429 & 14.3 & 17.7 & (13.5, 19.6) & 94.6 & 0.25 \\
Defense and National Security & 21 & 1{,}429 & 14.3 & 16.9 & (12.8, 19.0) & 92.5 & 0.20 \\
Information Integrity and Platforms & 18 & 1{,}428 & 14.0 & 15.7 & (13.1, 15.2) & 90.0 & 0.14 \\
Foreign Policy and Geopolitics & 22 & 1{,}429 & 13.5 & 14.7 & (10.1, 16.6) & 81.2 & ref \\
\bottomrule
\end{tabularx}
\endgroup
\apanote{Scenario-level adversarial effectivity on the 2,001-position scale, computed on the
independent scenario means after excluding the single-leaf macroeconomic domain. Columns:
number of directional opinion leaves, analysed scenarios, median (Mdn) and mean effectivity,
the interquartile range, the share of scenarios moving toward the attacker goal, and the
absolute rank-biserial effect of the scenario-level Mann-Whitney contrast against the
least-movable reference domain (foreign policy). Omnibus Kruskal-Wallis $H = 579.1$,
$p = 6.8\mathrm{e}{-123}$, $\varepsilon^2 = 0.067$.}
\end{table}
